% The stylized facts on the JumpHMM slide (spec, slide 4), after Fig. 3 of the JumpHMM paper
% (HMM-w-jumps-paper/arxiv-paper/figs/main/Fig03-Model-Comparison.pdf): real SPY daily growth 2014–2024 against 100
% generator paths of the same length. Input once in the preamble, after vnslides (for vnplotpanel, vnplotgrid and
% vnband).
%
% \vnstylizedfacts{<width>}{<height>} draws three panels side by side, each with a plot area <width> by <height>,
% from talk/data/sf_density.csv, sf_qq.csv and sf_abs_acf.csv (scripts/30_talk_data.jl):
%   (a) fat tails: kernel density of g (per 1/yr), with the normal of SPY's mean and SD dashed;
%   (b) tail quantiles: the generator's quantiles against SPY's, 0.1st to 99.9th percentile, on the 45° line;
%   (c) volatility clustering: ACF of |g| at lags 1–252 (trading days).
% The generator is its median with a 10th–90th percentile band (whiskers in (b)), and real SPY is the gold line, as SPY is on the
% dot plots.
% Plot areas share their bottom edge; each panel sits a fixed gap right of the last, room for its own y labels.
% Negative tick labels are written out with a true minus sign, and every tick label is in the text font.

% the look shared by the three panels
\pgfplotsset{
  vnsf/.style={
    scale only axis,
    axis background/.style={fill=vnplotpanel},
    axis line style={draw=none},
    grid=major, major grid style={draw=vnplotgrid, line width=0.6pt},
    tick style={draw=none},
    tick label style={font=\fontsize{7}{8.4}\selectfont, color=vndottick},
    xticklabel style={/pgf/number format/assume math mode=true},
    yticklabel style={/pgf/number format/assume math mode=true, /pgf/number format/fixed},
    label style={font=\fontsize{7.5}{9}\selectfont, color=vnslidebody},
    xlabel style={yshift=3pt},
    ylabel style={yshift=-3pt},
    % one box height for every title, so titles with and without descenders share a baseline
    title style={font=\fontsize{8.5}{10}\vnmedium, text=vnslideink, yshift=-3pt, text height=1.6ex,
      text depth=0.45ex},
  },
}

% \vnsfband[<median line width>]{<file>}{<x column>}: the generator's 10th–90th percentile band and its median line
\newcommand{\vnsfband}[3][0.9pt]{%
  \addplot[name path=vnsflo, draw=none, forget plot] table[col sep=comma, x=#3, y=syn_p10] {#2};
  \addplot[name path=vnsfhi, draw=none, forget plot] table[col sep=comma, x=#3, y=syn_p90] {#2};
  \addplot[fill=vnband, fill opacity=0.3, draw=none, forget plot] fill between[of=vnsflo and vnsfhi];
  \addplot[draw=vnslidelink, line width=#1, mark=none] table[col sep=comma, x=#3, y=syn_p50] {#2};
}

\newcommand{\vnstylizedfacts}[2]{%
  \begin{tikzpicture}
    % (a) the density, with the normal of SPY's mean and SD dashed; the generator's median sits on SPY's line, so it
    % is drawn wide underneath and SPY thin on top, and both show
    \begin{axis}[vnsf, name=vnsfa, width=#1, height=#2,
      title={fat tails},
      xmin=-10, xmax=10, xtick={-10,-5,0,5,10}, xticklabels={−10,−5,0,5,10},
      % no label at the bottom edge, 0, so the first x label has the corner to itself
      ymin=0, ymax=0.3, ytick={0.1,0.2,0.3},
      xlabel={daily growth, 1/yr}, ylabel={density},
    ]
      \vnsfband[2.4pt]{data/sf_density.csv}{x}
      \addplot[draw=vndottick, line width=0.8pt, dash pattern=on 2.5pt off 1.5pt, mark=none]
        table[col sep=comma, x=x, y=normal] {data/sf_density.csv};
      \addplot[draw=vnfigaccent, line width=0.9pt, line join=round, mark=none]
        table[col sep=comma, x=x, y=real] {data/sf_density.csv};
    \end{axis}
    % (b) the generator's quantiles against SPY's, on the 45° line
    \begin{axis}[vnsf, name=vnsfb, width=#1, height=#2,
      at={(vnsfa.south east)}, anchor=south west, xshift=1.25cm,
      title={tail quantiles},
      % labels inside the corners only, so the x and y labels never meet
      xmin=-20, xmax=20, xtick={-10,0,10}, xticklabels={−10,0,10},
      ymin=-20, ymax=20, ytick={-10,0,10}, yticklabels={−10,0,10},
      xlabel={SPY, 1/yr}, ylabel={generator, 1/yr},
    ]
      \addplot[draw=vndottick, line width=0.6pt, mark=none] coordinates {(-20,-20) (20,20)};
      \addplot[draw=vnslidelink, line width=0.7pt, mark=*, mark size=1.3pt, only marks,
        mark options={fill=vnslidelink, draw=vnslidelink},
        error bars/.cd, y dir=both, y explicit, error bar style={line width=0.7pt, draw=vnslidelink},
        error mark options={draw=none}]
        table[col sep=comma, x=real, y=syn_p50, y error plus expr={\thisrow{syn_p90} - \thisrow{syn_p50}},
          y error minus expr={\thisrow{syn_p50} - \thisrow{syn_p10}}] {data/sf_qq.csv};
    \end{axis}
    % (c) the ACF of |g| out to a year of lags
    \begin{axis}[vnsf, name=vnsfc, width=#1, height=#2,
      at={(vnsfb.south east)}, anchor=south west, xshift=1.25cm,
      title={volatility clustering},
      xmin=1, xmax=252, xtick={1,63,126,189,252},
      ymin=-0.05, ymax=0.45, ytick={0,0.2,0.4},
      xlabel={lag, trading days}, ylabel={ACF of $|g|$},
    ]
      \vnsfband{data/sf_abs_acf.csv}{lag}
      \addplot[draw=vnfigaccent, line width=1.0pt, line join=round, mark=none]
        table[col sep=comma, x=lag, y=real] {data/sf_abs_acf.csv};
    \end{axis}
  \end{tikzpicture}%
}
